%!TEX root = ../../main.tex
%% ===============================================================
%% 10.4 타당성 위협
%% 파일: chapters/10_discussion/04_threats.tex
%% 상위: chapters/10_discussion/chapter.tex
%% 정의 라벨: sec:disc-threats
%% 참조 라벨: tab:label-weights
%% 포함 객체: -
%% 인용: -
%% 상태: 초안 (docs/OUTLINE.md 참고)
%% ===============================================================

\section{타당성 위협}
\label{sec:disc-threats}

\paragraph{내적 타당성.} 무누출 계약은 단위 시험으로 강제되지만, Riot API 자체의 결함(예: 누락된 프레임, 지연된 사건 타임스탬프)은 통제할 수 없다. 교전의 \overlapPct{}\%는 서로 다른 장소에서 동시에 벌어져 시간적으로 겹치는데, 독립 표본을 가정하는 모델에는 무해하지만 경기 전체를 순차 처리하는 모델에는 라벨 누출이 될 수 있어 클리핑·마스킹이 필요하다.

\paragraph{외적 타당성.} 데이터는 한국 서버 상위 티어 랭크 경기이다. 다른 지역, 낮은 티어, 프로 경기로의 일반화는 검증되지 않았다. 세 패치에 걸친 시간순 홀드아웃은 단기 공변량 이동만을 다루며, 시즌 단위의 큰 변화(지도 개편, 신규 오브젝트)는 별도 검증이 필요하다.

\paragraph{구성 타당성.} 교환가치 라벨의 가중치(표 \ref{tab:label-weights})는 도메인 지식에 근거한 설계 선택이지 데이터로부터 추정된 값이 아니다. 가중치를 바꾸면 라벨이 바뀌므로 결과는 이 정의에 상대적이다. 검출기가 찾은 교전이 사람의 판단과 일치하는지는 진행 중인 주석 연구의 결과에 달려 있다.
